% =============================================================================
% 07a-algoritmo.tex - frame isolado, compartilhado entre a prelecao (07-sintese.tex) e o
% PDF avulso de bolso (isolados/algoritmo.tex). Fonte unica: editar aqui.
% =============================================================================
% ----------------------------------------------------------------- frame 39 --
\begin{frame}{\ftit{\faListOl}{O algoritmo}}
\CRar{6}
\begin{center}
\begin{tikzpicture}[node distance=0.16cm]
  \node[crnode={campoazul}{7.6cm}, font=\scriptsize, inner ysep=3.4pt] (a)
    {\textbf{1. Pergunte.} Use a frase de transi\c{c}\~ao, sem rodeio.};
  \node[crnode={campoazul}{7.6cm}, font=\scriptsize, inner ysep=3.4pt, below=of a] (b)
    {\textbf{2. Caracterize.} Frequ\^encia, dura\c{c}\~ao, intrusividade,
     capacidade de afastar, inten\c{c}\~ao, plano, prepara\c{c}\~ao, acesso.};
  \node[crnode={campoambar}{7.6cm}, font=\scriptsize, inner ysep=3.4pt, below=of b] (c)
    {\textbf{3. Identifique o que mudou.} \'Ultimas horas e \'ultimos dias.};
  \node[crnode={campoazul}{7.6cm}, font=\scriptsize, inner ysep=3.4pt, below=of c] (d)
    {\textbf{4. Formule.} Vulnerabilidades, modificadores, suicidalidade, acesso,
     capacidade de construir seguran\c{c}a.};
  \node[crnode={campobordo}{7.6cm}, font=\scriptsize, inner ysep=3.4pt, below=of d] (e)
    {\textbf{5. Decida o ambiente.} \'E poss\'ivel construir seguran\c{c}a fora
     daqui agora?};
  \node[crnode={campoverde}{7.6cm}, font=\scriptsize, inner ysep=3.4pt, below=of e] (f)
    {\textbf{6. Construa seguran\c{c}a.} Restri\c{c}\~ao de acesso pactuada e
     plano escrito com ele.};
  \node[crnode={campoverde}{7.6cm}, font=\scriptsize, inner ysep=3.4pt, below=of f] (g)
    {\textbf{7. Defina continuidade.} Data, local, respons\'avel, via de contato,
     conduta se piorar.};

  \foreach \i/\j/\c in {a/b/campoazul, b/c/campoambar, c/d/campoazul,
                        d/e/campobordo, e/f/campoverde, f/g/campoverde}
    \draw[crseta=\c] (\i.south) -- (\j.north);
\end{tikzpicture}
\end{center}
\end{frame}
